\chapter{Foot Fractures}

\section*{Definition and Overview}
A foot fracture is a break in one or more of the 26 bones of the foot, which include the toes (phalanges), the metatarsals of the midfoot, and the tarsal bones of the hindfoot, including the heel (calcaneus) and ankle region. Foot fractures are common injuries from falls, twisting, dropping heavy objects, and overuse, and they range from minor toe breaks that need little treatment to serious fractures that require surgery. The key distinction is whether the fracture is stable and well aligned, in which case it can heal in a boot or cast, or displaced or involving a joint, in which case surgery may be needed. This chapter explains the types of foot fractures, their symptoms, and their treatment.

\section*{Epidemiology}
Foot fractures are among the most common fractures seen in emergency departments, accounting for around 10% of all fractures. They are particularly common in sport, in the workplace, and in falls, and stress fractures affect runners, dancers, and military recruits. The metatarsals are the most commonly broken bones in the foot, and the fifth metatarsal has a distinctive fracture pattern (the Jones fracture) that heals poorly. Heel fractures, usually from falls from a height, are serious injuries with long recovery. Fractures are more common in older adults, in whom falls and osteoporosis combine.

\section*{Aetiology and Pathophysiology}
\begin{itemize}
    \item \textbf{Direct trauma:} dropping a heavy object on the foot, or kicking a hard object, breaking the toes and metatarsals
    \item \textbf{Twisting injuries:} inversion or forced rotation of the foot, fracturing the fifth metatarsal and tarsal bones
    \item \textbf{Falls from height:} landing on the feet fractures the heel bone, and often the spine too
    \item \textbf{Overuse:} repetitive loading causes stress fractures, most often in the metatarsals, in runners and people whose activity increases quickly
    \item \textbf{Osteoporosis:} weakened bone fractures with relatively minor force in older adults
    \item \textbf{Mechanism of injury matters:} the force determines the pattern, stability, and treatment
    \item \textbf{Classification:} fractures are described by which bone, whether displaced (ends not aligned), comminuted (shattered), open (through the skin), or intra-articular (into a joint)
    \item \textbf{Healing:} foot bones heal by forming new bone over 6--12 weeks, with protection from weight bearing during that time
\end{itemize}

\section*{Clinical Features}
\begin{itemize}
    \item Pain in the foot, immediate and made worse by weight bearing or pressure
    \item Swelling, bruising, and tenderness over the fracture site
    \item Difficulty walking or bearing weight
    \item Deformity, shortening, or abnormal movement of a toe
    \item In open fractures: a wound over the fracture, which is an emergency
    \item Numbness or tingling if nerves are affected
    \item Stress fractures: gradual onset pain with activity, easing with rest
    \item Ankle involvement: pain and swelling around the ankle joint with hindfoot fractures
\end{itemize}

\section*{Differential Diagnosis}
Foot pain after injury has several causes.
\begin{itemize}
    \item Sprains and ligament injuries, which are common in the ankle and midfoot
    \item Contusions and bruising without fracture
    \item Stress fractures, which may not appear on early X-rays
    \item Tendon injuries, including rupture of the Achilles tendon
    \item Gout, infection, and other causes of an acutely painful, swollen foot
    \item Metatarsalgia and plantar fasciitis, which cause gradual pain
    \item The X-ray and examination distinguish these
\end{itemize}

\section*{When to Seek Medical Care}
Anyone with foot pain after an injury that limits weight bearing, or with significant swelling or bruising, should have the foot assessed, with an X-ray where indicated.

\textbf{Call 000 (or local emergency services) for:}
\begin{itemize}
    \item An open fracture, with the bone or a wound over the fracture site
    \item A foot that is pale, cold, numb, or blue after injury
    \item Severe deformity or an obviously displaced fracture
    \item Falls from height with foot pain, which may also injure the spine
    \item Severe pain, or signs of infection in a diabetic person's foot
\end{itemize}

\section*{Diagnosis and Investigations}
\begin{itemize}
    \item \textbf{History and examination:} the mechanism of injury, pain, swelling, and the ability to bear weight
    \item \textbf{X-rays:} the standard investigation; multiple views are taken, and the whole foot and ankle are checked
    \item \textbf{CT scan:} for complex fractures of the heel and midfoot, to plan surgery
    \item \textbf{MRI or bone scan:} for suspected stress fractures not visible on X-ray
    \item \textbf{Assessment of the neurovascular status:} checking pulses and sensation, especially in displaced and open fractures
\end{itemize}

\section*{Management}
\begin{itemize}
    \item \textbf{Minor toe fractures:} buddy taping to the neighbouring toe and a stiff-soled shoe; most heal in 3--6 weeks
    \item \textbf{Metatarsal fractures:} a walking boot or cast for 4--8 weeks, with gradual return to activity
    \item \textbf{Stress fractures:} rest from the aggravating activity, with a return to training that is gradual
    \item \textbf{Fifth metatarsal (Jones) fractures:} often treated with a cast for 6--8 weeks, with surgery for delayed healing
    \item \textbf{Displaced fractures:} surgery to realign and fix the bones with plates, screws, or wires
    \item \textbf{Heel fractures:} significant injuries requiring a period without weight bearing, rehabilitation, and often surgery
    \item \textbf{Open fractures:} emergency surgery to clean the wound and stabilise the bone, with antibiotics
    \item \textbf{Aftercare:} elevation, ice, pain relief, and physiotherapy to restore strength and movement
\end{itemize}

\section*{Complications}
\begin{itemize}
    \item Delayed healing and non-union, especially in Jones fractures and in smokers
    \item Arthritis of the foot joints after fractures that enter a joint
    \item Chronic pain and stiffness
    \item Nerve damage and, in open or severe injuries, infection
    \item Compartment syndrome, a rare emergency with severe pain and swelling
    \item Long-term difficulty with footwear and activity after heel fractures
\end{itemize}

\section*{Prognosis and Outcomes}
Most foot fractures heal well with simple treatment, and minor toe and metatarsal fractures usually recover fully within a few months. Stress fractures settle with rest and graduated return to activity. The more serious fractures, particularly of the heel, can leave permanent stiffness and arthritis, though surgery and rehabilitation improve outcomes. In people with good bone health and non-smoking, healing is reliable. The key to a good outcome is an accurate diagnosis, appropriate immobilisation, and a structured return to activity.

\section*{Prevention and Self-Care}
\begin{itemize}
    \item Wear appropriate, well-fitting footwear for the activity
    \item Use protective footwear in hazardous work
    \item Build up running and exercise loads gradually to avoid stress fractures
    \item Keep bones strong with calcium, vitamin D, and weight-bearing exercise
    \item Treat osteoporosis and prevent falls in older adults
    \item Follow the weight-bearing instructions after a fracture
    \item Elevate and ice the foot after injury, and seek assessment for persistent pain
    \item Do not return to sport before the doctor or physiotherapist advises
\end{itemize}

\section*{Sources and Further Reading}
The following reputable sources provide further information for healthcare students and patients seeking reliable, current advice about foot fractures.
\begin{itemize}
    \item Healthdirect Australia. \textit{Foot fractures.} https://www.healthdirect.gov.au/
    \item Sports Medicine Australia. \textit{Foot injuries.} https://sma.org.au/
    \item Australian Orthopaedic Association. \textit{Fracture care.} https://www.aoa.org.au/
    \item NHS. \textit{Foot fracture.} https://www.nhs.uk/conditions/broken-toe/
    \item Mayo Clinic. \textit{Foot fracture.} https://www.mayoclinic.org/
    \item MedlinePlus. \textit{Foot fractures.} https://medlineplus.gov/
\end{itemize}
